% !TEX TS-program = xelatex
% !TEX encoding = UTF-8 Unicode
% !Mode:: "TeX:UTF-8"

\documentclass{resume}
\usepackage{zh_CN-Adobefonts_external} % Simplified Chinese Support using external fonts (./fonts/zh_CN-Adobe/)
%\usepackage{zh_CN-Adobefonts_internal} % Simplified Chinese Support using system fonts
\usepackage{linespacing_fix} % disable extra space before next section
\usepackage{cite}
\usepackage{hyperref}

% 字体设置：标题用黑体，正文用宋体（使用Adobe字体）
\usepackage{fontspec}
\setCJKmainfont[
  Path=fonts/zh_CN-Adobe/,
  Extension=.otf,
  BoldFont=AdobeHeitiStd-Regular.otf,
  ItalicFont=AdobeKaitiStd-Regular.otf,
  SmallCapsFont=AdobeHeitiStd-Regular.otf
]{AdobeSongStd-Light.otf}
\setCJKsansfont[
  Path=fonts/zh_CN-Adobe/,
  Extension=.otf
]{AdobeHeitiStd-Regular.otf}
\setCJKmonofont[
  Path=fonts/zh_CN-Adobe/,
  Extension=.otf
]{AdobeSongStd-Light.otf}
\begin{document}
\pagenumbering{gobble} % suppress displaying page number

\name{朱海}

% {E-mail}{mobilephone}{homepage}
% be careful of _ in emaill address
\contactInfo{Tel/VX:17816125272}{Email:1132863306@qq.com}{}{}
% {E-mail}{mobilephone}
% keep the last empty braces!
%\contactInfo{xxx@yuanbin.me}{(+86) 131-221-87xxx}{}

\vspace{0.5em}
\section{个人总结}
拥有2年阿里国际搜推算法经验，专注于推荐精排模型优化和生成式推荐。独立一作发表过AAAI、COLING等顶级会议论文，引用量90。对大模型/生成式推荐前沿技术保持热情。
\vspace{0.3em}

\section{工作经历}
\datedsubsection{\textbf{阿里国际AliExpress}, 推荐算法工程师}{2024.07-至今}
\begin{itemize}
\item 专注推荐精排模型商品建模方向，主导物流偏好建模，深度参与量价模型等工作，覆盖多个核心业务场景，服务亿级用户。
  \item 主导语义ID在排序模型中的探索与落地，设计并实现GateSID优化方案，提升新物品冷启动效率。
  \item 参与生成式排序Transformer模型（SORT）的设计与落地，实现订单\textcolor{red}{+6.35\%}、GMV\textcolor{red}{+5.47\%}的全场景业务提升。
\end{itemize}

\datedsubsection{\textbf{阿里国际Lazada}, 搜索算法工程师（实习）}{2023.07-2024.07}
\begin{itemize}
  \item 针对传统搜索相关性计算仅基于文本信息的局限，负责利用多模态大模型提升搜索相关性，提出Query-LIFE方法。
  \item 实现基于查询感知的图文模态融合，有效结合图像和标题信息，增强商品表示准确性。
  \item 线上取得订单UV提升\textcolor{red}{3.06\%}，GMV提升\textcolor{red}{3.19\%}，成果发表于COLING 2025。
\end{itemize}


\section{项目经历}

\datedsubsection{\textbf{语义ID精排建模（GateSID）}}{2025.07-2026.01}
\begin{itemize}
  \item \textit{情境}：推荐冷启动场景中，排序模型面临协同信号与语义信号权衡难题，新物品因缺乏协同信息难以获得足够曝光。
  \item \textit{任务}：设计能够根据物品成熟度动态平衡语义信息与协同信号的推荐框架，提升冷启动效率。
  \item \textit{行动}：提出GateSID框架，利用RQ-VAE将多模态特征离散化为分层语义ID，并构造用户语义ID行为序列；设计门控融合共享注意力机制，对冷启动物品依赖语义信号，对热门物品保留协同信号。
  \item \textit{结果}：GateSID在所有指标上均优于SOTA基线（COINS、SaviorRec），CTCVR AUC\textcolor{red}{+0.4\%}、GMV\textcolor{red}{+2.6\%}、订单量\textcolor{red}{+1.6\%},成果投稿于SIGIR 2026。
\end{itemize}

\datedsubsection{\textbf{工业级推荐系统排序Transformer模型（SORT）}}{2025.07-至今}
\begin{itemize}
  \item \textit{情境}：单场域数据独立训练导致用户全域意图建模缺失，传统排序模型scaling up收益递减。
  \item \textit{任务}：设计并落地基于生成式架构的transformer排序大模型，突破传统模型瓶颈。
  \item \textit{行动}：实现请求级样本组织、局部注意力；改进Transformer架构，引入特殊token、QKNorm、Attention Gate和Sparse MoE模块；优化训练/推理系统，联动工程团队开发稀疏注意力算子、算子融合，MFU从13\%提升至22\%。
  \item \textit{结果}：线上A/B测试订单\textcolor{red}{+6.35\%}、GMV\textcolor{red}{+5.47\%}；服务效率提升：延迟降低\textcolor{red}{44.67\%}，吞吐量\textcolor{red}{+121.33\%}；模型性能：CTR-AUC提升\textcolor{red}{2.1pt}，在数据/模型/序列长度维度均展现优秀扩展性。成果投稿KDD2026，详见\href{https://arxiv.org/pdf/2603.03988}{arXiv论文}
\end{itemize}

\datedsubsection{\textbf{量价模型优化}}{2024.07-2025.03}
\begin{itemize}
  \item \textit{情境}：原有排序模型对价格变动商品响应滞后，无法实时捕捉量价关系变化，影响流量分配效率。
  \item \textit{任务}：设计量价感知排序模型，实现对价格变化的快速响应，优化流量生态。
  \item \textit{行动}：深入研究量价关系，提出基于量价排序模型，增强价格信息表示，引入价格单调网络，优化CTR/CVR对价格变化的实时调整能力。
  \item \textit{结果}：降价商品流量\textcolor{red}{+41\%}，涨价商品流量\textcolor{red}{-18\%}，总订单\textcolor{red}{+1.9\%}，有效提升大盘效率，优化流量生态。
\end{itemize}

\datedsubsection{\textbf{物流偏好建模}}{2025.03-2025.09}
\begin{itemize}
  \item \textit{情境}：传统策略信号调整物流偏好方式僵硬且效率有损，无法精准识别优质物流服务。
  \item \textit{任务}：将物流履约信号融合进排序模型，实现自动化的物流质量识别与流量分配。
  \item \textit{行动}：构造用户/商品/交叉三侧物流信号，构建物流偏好建模模块，让模型自动识别并偏好优质物流服务。
  \item \textit{结果}：在GMV\textcolor{red}{+0.5\%至+1.1\%}前提下，物流优品订单\textcolor{red}{+1.0\%至+1.9\%}，物流差品订单\textcolor{red}{-1.5\%至+1.9\%}，签收时长降低\textcolor{red}{0.2至0.35天}。
\end{itemize}

\datedsubsection{\textbf{搜索相关性多模态（Query-LIFE）}}{2023.07-2024.07}
\begin{itemize}
  \item \textit{情境}：传统相关性计算方法仅基于产品标题和查询建模，无法充分利用图像信息，影响搜索准确性。
  \item \textit{任务}：设计多模态相关性计算方法，提升搜索结果的准确性和用户体验。
  \item \textit{行动}：提出Query-LIFE方法，采用基于查询的图文模态融合，有效结合图像和标题信息，利用查询感知的模态对齐增强商品表示准确性。
  \item \textit{结果}：线上A/B测试订单UV提升\textcolor{red}{3.06\%}，GMV提升\textcolor{red}{3.19\%}，成果发表于\href{https://aclanthology.org/2025.coling-industry.2.pdf}{COLING 2025}
\end{itemize}


%\datedsubsection{\textbf{对抗样本攻击课题}，独立一作}{2022.07-2022.11}
%\begin{itemize}
 % \item 针对基于深度学习的文本模型(TextCNN,LSTM,BERT)，利用了多种不同语义空间的替换词方法，提出BeamAttack攻击算法，在攻击成功率(ASR)，语义相似度(SIM)等指标上均超过之前的SOTA算法，该工作发表在\textbf{PAKDD2023(数据挖掘领域排名前五，录取率17.3\%)}。论文详见\url{https://link.springer.com/chapter/10.1007/978-3-031-33377-4_35}
%\end{itemize}

%\datedsubsection{\textbf{Transferable Adversarial Examples Can Efficiently Fool Topic Models
%}}{2021.8-2022.4}
%\begin{itemize}
%  \item 针对八个常用的主题模型(LDA,WTM等)，本人提出迁移攻击算法TopicAttack，评估攻击前后模型输出的KL散度差异作为攻击评价指标，生成的对抗样本对模型产生较大扰动。
%  \item 此外利用模型集成的方法，提出迁移性更高的攻击算法TopicAttack+，在三个数据集，八个主题模型上均取得比baseline更高的效果。这篇工作发表在\textbf{Computers \& Security（CCF-B）.}
%\end{itemize}

\section{技术影响}
\begin{itemize}[parsep=0.2ex]
  \item 在NLP/多模态/搜推方向均有相关经验，发表\textbf{AAAI、COLING}等多篇会议论文，引用量\textbf{90}，详见\href{https://scholar.google.com.hk/citations?hl=zh-CN&user=fKOITdsAAAAJ}{Google Scholar}。
  \item 对前沿技术保持好奇和关注，业余时间搭建AI论文阅读助手，包含前后端、数据库、大模型Agent等多种技术，提升算法同学阅读论文效率。累计收录\textbf{36w+}论文，触达\textbf{400+}用户，累计访问\textbf{3w+}，工作日DAU \textbf{30-40}。
\end{itemize}

\section{教育背景}
\datedsubsection{\textbf{中国科学技术大学},计算机技术,\textit{硕士}}{2021.09 - 2024.06}
\datedsubsection{\textbf{杭州电子科技大学},信息安全,\textit{学士}}{2017.09 - 2021.06}
\end{document}